Large language models encode numbers internally as base-10 digit representations, but how do they handle counting systems where the surface form encodes arithmetic rather than digits? French uses a vigesimal (base-20) system for numbers 70--99: ``quatre-vingt-dix-huit'' (98) literally means ``four-twenty-ten-eight'' ($4\times20+10+8$). We probe \mistral hidden representations to ask whether the model can decompose this arithmetic structure into standard base-10 digits. Using circular digit probes across four input formats---digit strings, English words, France French words, and Belgian French words (which use a simpler decimal system)---we find that French number word representations encode base-10 digits at 98.0\% accuracy, comparable to English (98.5\%) and Belgian French (99.0\%), and significantly higher than digit strings (91.0\%). The vigesimal structure causes only localized difficulty: the \emph{soixante-dix} (70s) range drops to 88.2\%, with errors revealing partial encoding of the ``soixante'' (60) component rather than the composed value. Our most striking finding is that number word representations peak at early layers (5--6), while digit string representations peak at the final layer (31), suggesting fundamentally different processing pathways for semantic and tokenized numerical inputs.
